\documentclass{lecture-kuznia}

\def\data{28.09.2026}
\def\place{Poręba Wielka}
\def\title{Nierówności, metoda stycznej}
\def\group{Finaliści}
\def\author{Autor: Stefan Świerczewski}
\def\lecturer{Prowadzący: Stefan Świerczewski}

\begin{document}
\begin{center}
	{\LARGE\textbf{\title{}}}
\end{center}

\section{Teoria}

\subsection{Po co nam styczna?}

Typowa nierówność ma postać
\[
	f(x_1)+f(x_2)+\dots+f(x_n)\geqslant C,
	\qquad x_1+\dots+x_n=S.
\]
Jeżeli spodziewamy się równości dla
\[
	x_1=\dots=x_n=x_0:=\frac Sn,
\]
to szukamy prostej $L(x)=\alpha x+\beta$, która dotyka wykresu $f$ w~punkcie
$x_0$. Warunki
\[
	L(x_0)=f(x_0),\qquad L'(x_0)=f'(x_0)
\]
dają jedyną kandydatkę:
\[
	L(x)=f(x_0)+f'(x_0)(x-x_0).
\]
Jeśli potrafimy wykazać, że $f(x)\geqslant L(x)$ w~całej \textbf{dopuszczalnej
dziedzinie}, to po zsumowaniu znika cała trudność:
\[
	\sum_{i=1}^n f(x_i)
	\geqslant
	\alpha\sum_{i=1}^n x_i+n\beta
	=\alpha S+n\beta=nf(x_0).
\]

\subsection{Styczna podpierająca i nierówność Jensena}

\begin{theorem}[prosta podpierająca]
	Niech $f$ będzie dwukrotnie różniczkowalna na przedziale $I$.
	Jeżeli $f''(x)\geqslant0$ dla $x\in I$, to dla dowolnego $x_0\in I$
	\[
		f(x)\geqslant f(x_0)+f'(x_0)(x-x_0)
		\qquad (x\in I).
	\]
	Dla $f''\leqslant0$ znak nierówności odwraca się.
\end{theorem}
\begin{proof}
	Załóżmy $f''\geqslant0$. Wtedy $f'$ jest niemalejąca. Dla $x>x_0$ mamy
	\[
		f(x)-f(x_0)=\int_{x_0}^{x}f'(t)\,dt
		\geqslant f'(x_0)(x-x_0).
	\]
	Dla $x<x_0$ całkujemy w~przeciwną stronę i~otrzymujemy tę samą nierówność.
	Przypadek wklęsły stosujemy do funkcji $-f$.
\end{proof}

\begin{theorem}[nierówność Jensena]
	Niech $f$ będzie różniczkowalną funkcją wypukłą na przedziale $I$, niech
	$x_1,\dots,x_n\in I$ oraz niech $\lambda_1,\dots,\lambda_n\geqslant0$
	spełniają $\lambda_1+\dots+\lambda_n=1$. Wtedy
	\[
		f\left(\lambda_1x_1+\dots+\lambda_nx_n\right)
		\leqslant
		\lambda_1f(x_1)+\dots+\lambda_nf(x_n).
	\]
	W~szczególności, dla równych wag,
	\[
		f\left(\frac{x_1+\dots+x_n}{n}\right)
		\leqslant\frac{f(x_1)+\dots+f(x_n)}{n}.
	\]
	Jeżeli $f$ jest wklęsła, oba znaki nierówności odwracają się.
\end{theorem}
\begin{proof}
	Oznaczmy $x_0=\lambda_1x_1+\dots+\lambda_nx_n$. Styczna podpierająca
	w~punkcie $x_0$ daje
	\[
		f(x_i)\geqslant f(x_0)+f'(x_0)(x_i-x_0).
	\]
	Mnożymy te nierówności przez $\lambda_i$ i~sumujemy. Składniki liniowe
	zerują się, ponieważ $\sum_i\lambda_i(x_i-x_0)=0$, więc otrzymujemy
	$\sum_i\lambda_if(x_i)\geqslant f(x_0)$.
\end{proof}

Jensen jest więc zsumowaną nierównością dla stycznych podpierających.

\begin{example}\label{E1}
	Niech $x_1,\dots,x_n>0$ oraz $x_1+\dots+x_n=n$. Udowodnij, że
	\[
		\frac1{x_1}+\dots+\frac1{x_n}\geqslant n.
	\]
\end{example}
\begin{solution*}[\ref{E1}]
	Funkcja $f(x)=1/x$ jest wypukła, a~jej styczna w~punkcie $x_0=1$ ma
	równanie $L(x)=2-x$. Zatem
	\[
		\frac1x\geqslant2-x,
	\]
	co jest równoważne nierówności $(x-1)^2/x\geqslant0$. Po zsumowaniu
	dostajemy tezę.
\end{solution*}

\subsection{Jak znaleźć punkt styczności}

W~praktyce działamy w~następującej kolejności:
\begin{enumerate}
	\item normalizujemy warunek, najczęściej do $x_1+\dots+x_n=n$ albo $1$;
	\item zgadujemy przypadek równości;
	\item grupujemy wyrażenie jako sumę wartości jednej funkcji;
	\item liczymy styczną w~punkcie równości;
	\item sprawdzamy znak różnicy, \textbf{pamiętając o~dziedzinie};
	\item dopiero na końcu sumujemy oszacowania.
\end{enumerate}

\subsection{Jak zgadnąć styczną bez pochodnej}

Niech $x_0$ będzie przewidywanym punktem równości i~niech
$L(x)=Ax+B$. Jeżeli $L$ jest prostą podpierającą, to różnica
$f(x)-L(x)$ zeruje się w~$x_0$, ale nie zmienia tam znaku. Po
sprowadzeniu tej różnicy do wspólnego mianownika (o~stałym znaku)
jej licznik powinien więc mieć pierwiastek co najmniej podwójny, czyli
czynnik $(x-x_0)^2$. Jest to algebraiczny odpowiednik warunków
$L(x_0)=f(x_0)$ oraz $L'(x_0)=f'(x_0)$.

Na przykład dla $f(x)=1/x$ i~$x_0=1$ szukamy $L(x)=A+Bx$ z~tożsamości
\[
	1-x(A+Bx)=C(x-1)^2.
\]
Porównanie współczynników daje $C=1$, $A=2$, $B=-1$, a~więc
\[
	L(x)=2-x,\qquad
	\frac1x-(2-x)=\frac{(x-1)^2}{x}\geqslant0.
\]
Podobnie, aby znaleźć górną styczną do $f(x)=x/(1+x)$ w~punkcie $1$,
układamy
\[
	(Ax+B)(1+x)-x=C(x-1)^2.
\]
Stąd $A=B=C=1/4$, zatem $L(x)=(x+1)/4$.

Wymuszenie kwadratu znajduje jedynie \emph{kandydatkę} na styczną.
Trzeba jeszcze sprawdzić znak całej różnicy w~pełnej dziedzinie zadania.
Dla punktu brzegowego albo funkcji niegładkiej argument z~czynnikiem
kwadratowym nie musi działać.

\newpage
\section{Zadania}

\begin{problem}\label{Z1}
	Niech $x_1,\dots,x_n\geqslant0$ oraz $x_1+\dots+x_n=n$. Wykaż, że
	\[
		\sum_{i=1}^n\frac{x_i}{1+x_i}\leqslant\frac n2.
	\]
	Wskaż przypadek równości.
\end{problem}

\begin{problem}\label{Z2}
	Niech $x_1,\dots,x_n>0$ oraz $x_1+\dots+x_n=n$. Udowodnij, że
	\[
		\sum_{i=1}^n\frac1{x_i^2+x_i+1}\geqslant\frac n3.
	\]
\end{problem}

\begin{problem}\label{ZKwadratWazony}
	Niech $a,b,c>0$. Udowodnij, że
	\[
		a(a+b-c)^2+b(b+c-a)^2+c(c+a-b)^2
		\geqslant\frac{(a^2+b^2+c^2)^2}{a+b+c}.
	\]
	Wyznacz przypadek równości.
\end{problem}

\begin{problem}\label{Z6}
	Udowodnij, że dla każdej liczby naturalnej $n\geqslant2$ zachodzi
	\[
		n!<\left(\frac{n+1}{2}\right)^n.
	\]
\end{problem}

\begin{problem}\label{Z8}
	Niech $x_1,\dots,x_n>0$ oraz $x_1+\dots+x_n=1$. Udowodnij, że
	\[
		\prod_{i=1}^n\left(1+\frac1{x_i}\right)\geqslant(n+1)^n.
	\]
\end{problem}

\begin{problem}\label{Z9}
	Niech $a,b,c>0$. Udowodnij, że
	\[
		\left(\frac{2a}{b+c}\right)^{\!2/3}
		+\left(\frac{2b}{c+a}\right)^{\!2/3}
		+\left(\frac{2c}{a+b}\right)^{\!2/3}\geqslant3.
	\]
\end{problem}

\begin{problem}\label{ZIloczynyJensen}
	Niech $n\geqslant2$ oraz $a_1,\dots,a_n\in(0,1/2)$. Udowodnij, że
	\[
		\frac{\prod_{i=1}^n a_i}
		{\left(\sum_{i=1}^n a_i\right)^n}
		\leqslant
		\frac{\prod_{i=1}^n(1-a_i)}
		{\left(\sum_{i=1}^n(1-a_i)\right)^n}.
	\]
	Wyznacz przypadek równości.
\end{problem}

\begin{problem}\label{ZE2}
	Niech $a,b,c\geqslant-\frac34$ oraz $a+b+c=1$. Udowodnij, że
	\[
		\frac{a}{a^2+1}+\frac{b}{b^2+1}+\frac{c}{c^2+1}
		\leqslant\frac9{10}.
	\]
\end{problem}

\begin{problem}\label{ZKatyCosinusy}
	Niech $\alpha,\beta,\gamma$ będą miarami kątów trójkąta. Udowodnij, że
	\[
		\cos\alpha+\cos\beta+\cos\gamma
		\geqslant12\cos\alpha\cos\beta\cos\gamma.
	\]
	Wyznacz przypadek równości.
\end{problem}

\begin{problem}\label{Z11}
	Niech $x_1,\dots,x_n>1$. Udowodnij, że
	\[
		\sum_{i=1}^n x_i+\frac{n}{\sqrt[n]{x_1\cdots x_n}}
		\geqslant
		\sum_{i=1}^n\frac1{x_i}+n\sqrt[n]{x_1\cdots x_n}.
	\]
\end{problem}

\begin{problem}\label{Z10}
	Niech $a,b,c>0$ oraz $a^3+b^3+c^3=3$. Udowodnij, że
	\[
		\frac1{3-2a}+\frac1{3-2b}+\frac1{3-2c}\geqslant3.
	\]
\end{problem}

\begin{problem}\label{Z14}
	Niech $a,b,c>0$. Udowodnij, że
	\[
		\frac{(b+c-a)^2}{(b+c)^2+a^2}
		+\frac{(c+a-b)^2}{(c+a)^2+b^2}
		+\frac{(a+b-c)^2}{(a+b)^2+c^2}
		\geqslant\frac35.
	\]
\end{problem}

\begin{problem}\label{ZRaviJensen}
	Niech $a,b,c$ będą długościami boków trójkąta. Udowodnij, że
	\[
		a^2b(a-b)+b^2c(b-c)+c^2a(c-a)\geqslant0.
	\]
	Wyznacz przypadek równości.
\end{problem}

\begin{problem}\label{Z12}
	Niech $a,b,c>0$ oraz
	\[
		57a+88b+125c\geqslant1148.
	\]
	Wyznacz najmniejszą możliwą wartość wyrażenia
	\[
		a^3+b^3+c^3+5a^2+5b^2+5c^2.
	\]
\end{problem}

\begin{problem}\label{Z13}
	Niech $a,b,c>0$ oraz $a+b+c=3$. Udowodnij, że
	\[
		18\sum_{\rm cyc}\frac1{(3-a)(4-a)}+2(ab+bc+ca)\geqslant15.
	\]
\end{problem}

\begin{problem}\label{ZHomoJensen}
	Niech $a,b,c>0$ oraz
	\[
		\frac1a+\frac1b+\frac1c=a+b+c.
	\]
	Udowodnij, że
	\[
		\frac1{(2a+b+c)^2}+\frac1{(a+2b+c)^2}
		+\frac1{(a+b+2c)^2}\leqslant\frac3{16}.
	\]
\end{problem}

\begin{problem}\label{Z15}
	Niech $n\geqslant4$, $a_1,\dots,a_n\geqslant0$ oraz
	$a_1+\dots+a_n=2$. Wyznacz najmniejszą możliwą wartość wyrażenia
	\[
		\frac{a_1}{1+a_2^2}+\frac{a_2}{1+a_3^2}
		+\dots+\frac{a_n}{1+a_1^2}.
	\]
\end{problem}

\ifshowsolutions
\newpage
\section{Wskazówki}

\begin{hint*}[\ref{Z1}]
	Porównaj $f(x)=x/(1+x)$ z~jej styczną w~punkcie $1$.
\end{hint*}

\begin{hint*}[\ref{Z2}]
	Punkt równości to $1$. Odejmij od funkcji $1/(x^2+x+1)$ jej styczną
	w~tym punkcie i~rozłóż różnicę na czynniki.
\end{hint*}

\begin{hint*}[\ref{ZKwadratWazony}]
	Zastosuj ważoną nierówność Jensena do funkcji $f(x)=x^2$, przyjmując
	wagi proporcjonalne do $a,b,c$.
\end{hint*}

\begin{hint*}[\ref{Z6}]
	Zastosuj ścisłą wklęsłość $\log t$ do liczb $1,2,\dots,n$.
\end{hint*}

\begin{hint*}[\ref{Z8}]
	Zastosuj Jensena do wypukłej funkcji
	$f(t)=\log(1+1/t)$.
\end{hint*}

\begin{hint*}[\ref{Z9}]
	Po normalizacji $a+b+c=1$ wykaż punktowo
	$\left(2t/(1-t)\right)^{2/3}\geqslant3t$.
\end{hint*}

\begin{hint*}[\ref{ZIloczynyJensen}]
	Zastosuj Jensena do funkcji
	$f(x)=\log\frac{x}{1-x}$ na przedziale $(0,1/2)$.
\end{hint*}

\begin{hint*}[\ref{ZE2}]
	Dla $f(x)=x/(x^2+1)$ zapisz styczną w~punkcie $1/3$ i~sprawdź jej
	położenie względem wykresu na całym przedziale $[-3/4,\infty)$.
\end{hint*}

\begin{hint*}[\ref{ZKatyCosinusy}]
	Dla trójkąta ostrokątnego najpierw ogranicz
	$\cos\alpha\cos\beta\cos\gamma$ za pomocą wklęsłości funkcji
	$\log(\cos x)$, a~następnie użyj AM--GM.
\end{hint*}

\begin{hint*}[\ref{Z11}]
	Podstaw $u_i=\log x_i$ i~zastosuj Jensena do funkcji
	$f(u)=e^u-e^{-u}$ na $(0,\infty)$.
\end{hint*}

\begin{hint*}[\ref{Z10}]
	Najpierw sprawdź, że wszystkie mianowniki są dodatnie. Podstaw $u=x^3$
	i~porównaj funkcję $1/(3-2\sqrt[3]{u})$ z~jej styczną w~punkcie $u=1$.
\end{hint*}

\begin{hint*}[\ref{Z14}]
	Znormalizuj $a+b+c=1$ i~zapisz każdy składnik w~postaci
	$2-1/(2a^2-2a+1)$. Potem użyj stycznej w~punkcie $1/3$.
\end{hint*}

\begin{hint*}[\ref{ZRaviJensen}]
	Zastosuj podstawienie
	$a=y+z$, $b=z+x$, $c=x+y$. Po uproszczeniu użyj ważonego Jensena
	dla funkcji $f(t)=t^2$.
\end{hint*}

\begin{hint*}[\ref{Z12}]
	Dla $f(x)=x^3+5x^2$ zapisz styczne w~punktach $3$, $4$ i~$5$.
	Ich współczynniki kierunkowe występują wprost w~założeniu.
\end{hint*}

\begin{hint*}[\ref{Z13}]
	Zastąp $2(ab+bc+ca)$ wyrażeniem z~$a^2+b^2+c^2$, a~następnie
	zastosuj styczną w~punkcie $1$ do otrzymanej funkcji jednej zmiennej.
\end{hint*}

\begin{hint*}[\ref{ZHomoJensen}]
	Oznacz $s=a+b+c$ i~najpierw ujednorodnij nierówność, mnożąc jej
	lewą stronę przez $s^2$, a~prawą zastępując przez
	$\frac3{16}s(1/a+1/b+1/c)$. Potem przyjmij $s=3$ i~zastosuj Jensena
	do odpowiedniej funkcji jednej zmiennej.
\end{hint*}

\begin{hint*}[\ref{Z15}]
	Najpierw wykaż dla $x\geqslant0$ oszacowanie
	\[
		\frac1{1+x^2}\geqslant1-\frac x2.
	\]
	Następnie udowodnij, że
	$\sum_{i=1}^n a_ia_{i+1}\leqslant1$, gdzie indeksy rozumiemy cyklicznie.
\end{hint*}

\newpage
\section{Rozwiązania}

\begin{solution*}[\ref{Z1}]
	Dla $f(x)=x/(1+x)$ styczna w~punkcie $1$ ma postać
	$L(x)=(x+1)/4$. Różnica
	\[
		L(x)-f(x)=\frac{(x-1)^2}{4(1+x)}\geqslant0.
	\]
	Sumując, dostajemy
	\[
		\sum_{i=1}^n\frac{x_i}{1+x_i}
		\leqslant\frac14\sum_{i=1}^n(x_i+1)=\frac n2.
	\]
	Równość zachodzi dokładnie dla $x_1=\dots=x_n=1$.
\end{solution*}

\begin{solution*}[\ref{Z2}]
	Punkt równości to $x=1$. Styczna do $f(x)=1/(x^2+x+1)$ w~tym punkcie
	ma równanie $L(x)=(2-x)/3$. Kluczowa faktoryzacja brzmi
	\[
		\frac1{x^2+x+1}-\frac{2-x}{3}
		=\frac{(x-1)^2(x+1)}{3(x^2+x+1)}\geqslant0.
	\]
	Po zsumowaniu prawych stron otrzymujemy
	$\frac13(2n-\sum x_i)=n/3$.
\end{solution*}

\begin{solution*}[\ref{ZKwadratWazony}]
	Oznaczmy $s=a+b+c$. Z~ważonej nierówności Jensena dla funkcji
	$f(x)=x^2$ otrzymujemy
	\[
		\frac{a}{s}(a+b-c)^2+\frac{b}{s}(b+c-a)^2
		+\frac{c}{s}(c+a-b)^2
		\geqslant
		\left(
			\frac{a(a+b-c)+b(b+c-a)+c(c+a-b)}{s}
		\right)^2.
	\]
	Wyrażenie w~liczniku po prawej stronie upraszcza się do
	$a^2+b^2+c^2$. Po pomnożeniu przez $s$ dostajemy tezę.
	Równość zachodzi wtedy i~tylko wtedy, gdy
	\[
		a+b-c=b+c-a=c+a-b,
	\]
	czyli gdy $a=b=c$.
\end{solution*}

\begin{solution*}[\ref{Z6}]
	Funkcja $\log t$ jest ściśle wklęsła. Jensen zastosowany do liczb
	$1,2,\dots,n$ daje
	\[
		\frac1n\sum_{k=1}^n\log k
		<\log\left(\frac1n\sum_{k=1}^n k\right)
		=\log\frac{n+1}{2}.
	\]
	Nierówność jest ostra, bo dla $n\geqslant2$ liczby nie są wszystkie równe.
	Po eksponencjacji otrzymujemy tezę.
\end{solution*}

\begin{solution*}[\ref{Z8}]
	Dla $t>0$ funkcja
	\[
		f(t)=\log\left(1+\frac1t\right)
	\]
	jest wypukła, ponieważ
	\[
		f''(t)=\frac{2t+1}{t^2(t+1)^2}>0.
	\]
	Z~nierówności Jensena
	\[
		\frac1n\sum_{i=1}^n\log\left(1+\frac1{x_i}\right)
		\geqslant
		\log\left(1+\frac1{(x_1+\dots+x_n)/n}\right)
		=\log(n+1).
	\]
	Eksponencjacja kończy dowód. Równość zachodzi dla
	$x_1=\dots=x_n=1/n$.
\end{solution*}

\begin{solution*}[\ref{Z9}]
	Oznaczmy $s=a+b+c$ oraz $x=a/s$, $y=b/s$, $z=c/s$. Wtedy
	$x+y+z=1$. Rozważmy
	\[
		f(t)=\left(\frac{2t}{1-t}\right)^{2/3}
		\qquad(0<t<1).
	\]
	Styczna w~punkcie $t_0=1/3$ ma równanie $L(t)=3t$. Jest ona prostą
	podpierającą, ponieważ nierówność $f(t)\geqslant3t$ jest równoważna
	\[
		4-27t(1-t)^2=(3t-1)^2(4-3t)\geqslant0.
	\]
	Zatem
	\[
		\sum_{\rm cyc}\left(\frac{2a}{b+c}\right)^{2/3}
		=f(x)+f(y)+f(z)\geqslant3(x+y+z)=3.
	\]
	Równość zachodzi dla $a=b=c$.
\end{solution*}

\begin{solution*}[\ref{ZIloczynyJensen}]
	Na przedziale $(0,1/2)$ funkcja
	\[
		f(x)=\log\frac{x}{1-x}
	\]
	jest ściśle wklęsła, ponieważ
	\[
		f''(x)=\frac{2x-1}{x^2(1-x)^2}<0.
	\]
	Oznaczmy $\overline a=(a_1+\dots+a_n)/n$. Z~nierówności Jensena
	\[
		\frac1n\sum_{i=1}^n\log\frac{a_i}{1-a_i}
		\leqslant
		\log\frac{\overline a}{1-\overline a}
		=
		\log\frac{\sum_{i=1}^n a_i}{\sum_{i=1}^n(1-a_i)}.
	\]
	Po pomnożeniu przez $n$, potęgowaniu i~przestawieniu czynników
	otrzymujemy dokładnie tezę. Ze ścisłej wklęsłości wynika, że równość
	zachodzi wtedy i~tylko wtedy, gdy $a_1=\dots=a_n$.
\end{solution*}

\begin{solution*}[\ref{ZE2}]
	Dla $f(x)=x/(x^2+1)$ równość sugeruje punkt $x_0=1/3$. Styczna ma
	równanie
	\[
		L(x)=\frac{18}{25}x+\frac3{50}.
	\]
	Nie musimy badać znaku $f''$. Bezpośrednio
	\[
		L(x)-f(x)
		=\frac{(3x-1)^2(4x+3)}{50(x^2+1)}\geqslant0
		\qquad\left(x\geqslant-\frac34\right).
	\]
	Po zsumowaniu trzech oszacowań otrzymujemy
	\[
		\sum_{\rm cyc}f(a)
		\leqslant\frac{18}{25}(a+b+c)+3\cdot\frac3{50}
		=\frac9{10}.
	\]
	Ograniczenie dziedziny jest dokładnie tym, co czyni styczną podpierającą.
\end{solution*}

\begin{solution*}[\ref{ZKatyCosinusy}]
	Jeżeli trójkąt nie jest ostrokątny, to prawa strona jest niedodatnia,
	natomiast
	\[
		\cos\alpha+\cos\beta+\cos\gamma
		=1+4\sin\frac\alpha2\sin\frac\beta2\sin\frac\gamma2>0.
	\]
	Możemy więc założyć, że wszystkie trzy kąty są ostre, i~oznaczyć
	$p=\cos\alpha\cos\beta\cos\gamma>0$. Funkcja
	$f(x)=\log(\cos x)$ jest ściśle wklęsła na $(0,\pi/2)$, zatem Jensen
	daje
	\[
		\frac13\log p
		\leqslant
		\log\left(\cos\frac{\alpha+\beta+\gamma}{3}\right)
		=\log\frac12.
	\]
	Stąd $p\leqslant1/8$. Z~nierówności AM--GM oraz
	$p^{2/3}\leqslant1/4$ otrzymujemy
	\[
		\cos\alpha+\cos\beta+\cos\gamma
		\geqslant3p^{1/3}
		\geqslant12p.
	\]
	Równość zachodzi dokładnie dla
	$\alpha=\beta=\gamma=\pi/3$.
\end{solution*}

\begin{solution*}[\ref{Z11}]
	Oznaczmy $u_i=\log x_i>0$ oraz
	\[
		\overline u=\frac{u_1+\dots+u_n}{n},
		\qquad G=e^{\overline u}=\sqrt[n]{x_1\cdots x_n}.
	\]
	Funkcja $f(u)=e^u-e^{-u}$ jest wypukła na $(0,\infty)$, gdyż
	$f''(u)=e^u-e^{-u}>0$. Zatem Jensen daje
	\[
		\sum_{i=1}^n\left(x_i-\frac1{x_i}\right)
		=\sum_{i=1}^n f(u_i)
		\geqslant nf(\overline u)
		=n\left(G-\frac1G\right).
	\]
	Po przeniesieniu składników otrzymujemy tezę.
\end{solution*}

\begin{solution*}[\ref{Z10}]
	Z~warunku mamy $a,b,c<\sqrt[3]{3}<3/2$, zatem wszystkie mianowniki
	w~tezie są dodatnie. Dla $0<x<3/2$ zachodzi oszacowanie
	\[
		\frac1{3-2x}\geqslant\frac13+\frac23x^3.
	\]
	Jest to styczna w~punkcie $u=1$ do funkcji
	$g(u)=1/(3-2\sqrt[3]{u})$. Potrzebny znak można też sprawdzić bez
	badania wypukłości:
	\[
		\frac1{3-2x}-\left(\frac13+\frac23x^3\right)
		=\frac{2x(x-1)^2(2x+1)}{3(3-2x)}\geqslant0.
	\]
	Stosując tę nierówność kolejno do $a,b,c$, otrzymujemy
	\[
		\frac1{3-2a}+\frac1{3-2b}+\frac1{3-2c}
		\geqslant1+\frac23(a^3+b^3+c^3)=3.
	\]
	Równość zachodzi dla $a=b=c=1$.
\end{solution*}

\begin{solution*}[\ref{Z14}]
	Nierówność jest jednorodna, więc zakładamy $a+b+c=1$. Wtedy pierwszy
	składnik ma postać
	\[
		\frac{(1-2a)^2}{(1-a)^2+a^2}
		=2-\frac1{2a^2-2a+1}.
	\]
	Wystarczy zatem wykazać
	\[
		\sum_{\rm cyc}\frac1{2a^2-2a+1}\leqslant\frac{27}{5}.
	\]
	Dla $f(x)=1/(2x^2-2x+1)$ styczna w~punkcie $x_0=1/3$ ma równanie
	\[
		L(x)=\frac{54}{25}x+\frac{27}{25}.
	\]
	Co ważniejsze, dla $x>0$
	\[
		f(x)-L(x)
		=-\frac{2(3x-1)^2(6x+1)}{25(2x^2-2x+1)}\leqslant0.
	\]
	Sumujemy oszacowania:
	\[
		\sum_{\rm cyc}f(a)
		\leqslant\frac{54}{25}(a+b+c)+3\cdot\frac{27}{25}
		=\frac{27}{5}.
	\]
	Po cofnięciu przekształcenia otrzymujemy tezę. Równość zachodzi dla
	$a=b=c$.
\end{solution*}

\begin{solution*}[\ref{ZRaviJensen}]
	Z~nierówności trójkąta istnieją liczby $x,y,z>0$ takie, że
	\[
		a=y+z,\qquad b=z+x,\qquad c=x+y.
	\]
	Po podstawieniu i~uproszczeniu teza jest równoważna nierówności
	\[
		x^3z+y^3x+z^3y\geqslant xyz(x+y+z),
	\]
	a~po podzieleniu przez $xyz>0$ nierówności
	\[
		\frac{x^2}{y}+\frac{y^2}{z}+\frac{z^2}{x}\geqslant x+y+z.
	\]
	Stosujemy ważonego Jensena do wypukłej funkcji $f(t)=t^2$, z~wagami
	$y,z,x$:
	\begin{align*}
		yf\left(\frac{x}{y}\right)
		+zf\left(\frac{y}{z}\right)
		+xf\left(\frac{z}{x}\right)
		&\geqslant(x+y+z)
		f\left(\frac{x+y+z}{x+y+z}\right)\\
		&=x+y+z.
	\end{align*}
	Równość zachodzi wtedy i~tylko wtedy, gdy
	$x/y=y/z=z/x=1$, czyli $x=y=z$, co jest równoważne warunkowi
	$a=b=c$.
\end{solution*}

\begin{solution*}[\ref{Z12}]
	Rozważmy wypukłą na $(0,\infty)$ funkcję $f(x)=x^3+5x^2$.
	Jej pochodna ma postać $f'(x)=3x^2+10x$. Styczne odpowiednio
	w~punktach $3$, $4$ i~$5$ dają oszacowania
	\[
		f(a)\geqslant57a-99,\qquad
		f(b)\geqslant88b-208,\qquad
		f(c)\geqslant125c-375.
	\]
	Po zsumowaniu i~wykorzystaniu założenia otrzymujemy
	\[
		f(a)+f(b)+f(c)
		\geqslant57a+88b+125c-682
		\geqslant466.
	\]
	Dla $(a,b,c)=(3,4,5)$ założenie zachodzi z~równością, a~badane
	wyrażenie jest równe $72+144+250=466$. Szukane minimum wynosi więc
	$466$.
\end{solution*}

\begin{solution*}[\ref{Z13}]
	Ponieważ $a^2+b^2+c^2=9-2(ab+bc+ca)$, lewa strona jest równa
	\[
		9+\sum_{\rm cyc}\left(\frac{18}{(3-a)(4-a)}-a^2\right).
	\]
	Styczna w~punkcie $x_0=1$ do funkcji
	$f(x)=18/((3-x)(4-x))-x^2$ ma równanie $L(x)=(x+3)/2$. Dla
	$0<x<3$ mamy dokładnie
	\[
		f(x)-L(x)
		=\frac{x(x-1)^2(9-2x)}{2(3-x)(4-x)}\geqslant0.
	\]
	Sumując, otrzymujemy $9+\frac12((a+b+c)+9)=15$.
	Równość zachodzi dla $a=b=c=1$.
\end{solution*}

\begin{solution*}[\ref{ZHomoJensen}]
	Oznaczmy $s=a+b+c$. Korzystając z~założenia, możemy równoważnie
	udowodnić jednorodną nierówność
	\[
		\sum_{\rm cyc}\frac{s^2}{(s+a)^2}
		\leqslant\frac3{16}s\left(\frac1a+\frac1b+\frac1c\right).
	\]
	Możemy więc przyjąć $s=3$. Po podzieleniu obu stron przez $9$
	wystarczy wykazać
	\[
		\sum_{\rm cyc}\left(\frac{16}{(a+3)^2}-\frac1a\right)\leqslant0,
		\qquad a+b+c=3.
	\]
	Rozważmy funkcję
	\[
		f(x)=\frac{16}{(x+3)^2}-\frac1x
		\qquad(0<x\leqslant3).
	\]
	Mamy $f(1)=0$ oraz
	\[
		f''(x)=\frac{96}{(x+3)^4}-\frac2{x^3}\leqslant0.
	\]
	Ostatnia nierówność jest równoważna
	\[
		(x+3)^4-48x^3
		=(x-3)(x^3-33x^2-45x-27)\geqslant0,
	\]
	co zachodzi dla $0<x\leqslant3$. Z~wklęsłości funkcji $f$ i~nierówności
	Jensena otrzymujemy
	\[
		f(a)+f(b)+f(c)
		\leqslant3f\left(\frac{a+b+c}{3}\right)=3f(1)=0.
	\]
	Równość w~wyjściowym zadaniu zachodzi dla $a=b=c=1$.
\end{solution*}

\begin{solution*}[\ref{Z15}]
	Dla każdego $x\geqslant0$ zachodzi
	\[
		\frac1{1+x^2}\geqslant1-\frac x2,
	\]
	ponieważ
	\[
		\frac1{1+x^2}-\left(1-\frac x2\right)
		=\frac{x(x-1)^2}{2(1+x^2)}\geqslant0.
	\]
	Przyjmijmy $a_{n+1}=a_1$ i~oznaczmy
	$Q=\sum_{i=1}^n a_ia_{i+1}$. Pokażemy najpierw, że $Q\leqslant1$.
	Na zwartym zbiorze $a_i\geqslant0$, $\sum a_i=2$ funkcja $Q$ osiąga
	maksimum. Spośród punktów maksymalizujących wybierzmy taki, który ma
	najmniej dodatnich współrzędnych.

	Gdyby dwie dodatnie współrzędne $a_p,a_q$ nie sąsiadowały ze sobą
	w~cyklu, moglibyśmy pozostawić ich sumę bez zmian i~przenieść ją w~całości
	do jednej z~nich. Wyrażenie $Q$ jest przy takim przesunięciu funkcją
	liniową, więc jeden z~końców przedziału nie zmniejsza jego wartości.
	Otrzymalibyśmy maksimum z~mniejszą liczbą dodatnich współrzędnych,
	sprzecznie z~wyborem. Wszystkie dodatnie współrzędne muszą więc parami
	sąsiadować. W~cyklu długości co najmniej $4$ mogą być najwyżej dwie.
	Jeśli mają wartości $x,y$, to
	\[
		Q=xy\leqslant\left(\frac{x+y}{2}\right)^2=1.
	\]
	Korzystając teraz z~oszacowania jednowymiarowego, dostajemy
	\[
		\sum_{i=1}^n\frac{a_i}{1+a_{i+1}^2}
		\geqslant\sum_{i=1}^na_i\left(1-\frac{a_{i+1}}2\right)
		=2-\frac Q2\geqslant\frac32.
	\]
	Dla $(a_1,a_2,a_3,\dots,a_n)=(1,1,0,\dots,0)$ otrzymujemy wartość
	$3/2$, zatem jest ona szukanym minimum.
\end{solution*}

\fi

\end{document}
